% SHARD-ID: AQM-62-TANGENT-COTANGENT-WEYL-KINEMATICS-MORE-EXAMPLES
% SHARD-TITLE: More Tangent-Cotangent Weyl Kinematical Examples
% SHARD-SUMMARY: Continues AQM-61 with affine-line, affine-plane, crossing, and circle examples.
% SHARD-SUMMARY: Describes the actual Hilbert spaces for rational supports and finite-residue closed points.
% SHARD-SUMMARY: Confirms again that tangent-cotangent Weyl data are kinematics, not dynamics.
% SHARD-KEYWORDS: tangent, cotangent, Weyl, Hilbert space, F3, F5, affine line, affine plane, crossing, circle

\section{More Tangent-Cotangent Weyl Kinematical Examples}
\label{sec:tangent-cotangent-weyl-kinematics-more-examples}

\subsection{Status and Source Anchors}

This shard continues AQM-61 and uses the same local sources.  The operative
recipe is:
\[
  T_P^*X=\Omega^1_{A/k}\otimes_A\kappa(P),\qquad
  T_PX=\operatorname{Hom}_{\kappa(P)}(T_P^*X,\kappa(P)),
\]
\[
  E_P^{\rm geom}=T_PX\oplus T_P^*X,\qquad
  \mathcal H_P^{\rm geom}=\ell^2(T_PX).
\]
For a finite rational support \(S\),
\[
  \mathcal H_S^{\rm geom}=\bigotimes_{P\in S}\mathcal H_P^{\rm geom}.
\]

\subsection{Examples 6--10}

\paragraph{6. The affine line over \(\mathbb F_3\).}
Let \(X=\Spec k[x]\), \(k=\mathbb F_3\).  At the rational point \(P_a=(x-a)\),
\[
  T_{P_a}^*X=k\,dx,\qquad T_{P_a}X=k\,\partial_x.
\]
Thus each rational point contributes
\[
  \mathcal H_{P_a}^{\rm geom}=\ell^2(k)\cong\mathbb C^3.
\]
For the rational support \(S=\mathbb A^1(k)=\{0,1,2\}\),
\[
  \mathcal H_S^{\rm geom}\cong(\mathbb C^3)^{\otimes3},
  \qquad \dim\mathcal H_S^{\rm geom}=3^3.
\]
For a closed point of degree \(d\), \(K=\mathbb F_{3^d}\), \(T_PX=K\), and
\[
  \mathcal H_P^{\rm geom}=\ell^2(\mathbb F_{3^d})
  \cong\mathbb C^{3^d}.
\]

\paragraph{7. The affine plane over \(\mathbb F_3\).}
Let \(X=\Spec k[x,y]\), \(k=\mathbb F_3\).  At any rational point \(P=(a,b)\),
\[
  T_P^*X=k\,dx\oplus k\,dy,\qquad
  T_PX=k\,\partial_x\oplus k\,\partial_y.
\]
Hence
\[
  \mathcal H_P^{\rm geom}=\ell^2(k^2)\cong\mathbb C^9.
\]
For the rational support \(S=k^2\), there are nine sites, each with Hilbert
space \(\mathbb C^9\), so
\[
  \mathcal H_S^{\rm geom}\cong(\mathbb C^9)^{\otimes9},
  \qquad \dim\mathcal H_S^{\rm geom}=9^9=3^{18}.
\]

\paragraph{8. The crossing \(xy=0\) over \(\mathbb F_3\).}
Let
\[
  A=k[x,y]/(xy),\qquad k=\mathbb F_3.
\]
At a rational point \(P=(a,b)\) on \(xy=0\), the tangent equation is
\[
  b\,u+a\,v=0.
\]
At the origin this imposes no equation, so
\[
  T_{(0,0)}X=k^2,\qquad
  \mathcal H_{(0,0)}^{\rm geom}\cong\mathbb C^9.
\]
At the four other rational points
\[
  (1,0),(2,0),(0,1),(0,2)
\]
the tangent space is one-dimensional, so each contributes \(\mathbb C^3\).
For all rational points,
\[
  \mathcal H_{X(k)}^{\rm geom}\cong
  \mathbb C^9\otimes(\mathbb C^3)^{\otimes4},
  \qquad \dim\mathcal H_{X(k)}^{\rm geom}=3^6.
\]

\paragraph{9. The circle over \(\mathbb F_3\).}
Let
\[
  C=\Spec k[x,y]/(x^2+y^2-1),\qquad k=\mathbb F_3.
\]
The rational points are
\[
  (1,0),(2,0),(0,1),(0,2).
\]
At \(P=(a,b)\), the tangent equation is
\[
  a\,u+b\,v=0,
\]
because \(2\) is invertible in \(k\).  Each point is smooth and has
\(\dim_kT_PC=1\).  Therefore each site has
\[
  \mathcal H_P^{\rm geom}\cong\mathbb C^3.
\]
For the rational support,
\[
  \mathcal H_{C(k)}^{\rm geom}\cong(\mathbb C^3)^{\otimes4},
  \qquad \dim\mathcal H_{C(k)}^{\rm geom}=3^4.
\]

\paragraph{10. The circle over \(\mathbb F_5\).}
Let
\[
  C=\Spec \mathbb F_5[x,y]/(x^2+y^2-1).
\]
Since \(-1\) is a square in \(\mathbb F_5\), AQM-59 gives four rational points:
\[
  (1,0),(4,0),(0,1),(0,4).
\]
At \(P=(a,b)\), the tangent equation is again
\[
  a\,u+b\,v=0,
\]
so every rational tangent space is one-dimensional over \(\mathbb F_5\).
Consequently
\[
  \mathcal H_P^{\rm geom}=\ell^2(\mathbb F_5)\cong\mathbb C^5
\]
at each rational point, and
\[
  \mathcal H_{C(\mathbb F_5)}^{\rm geom}
  \cong(\mathbb C^5)^{\otimes4},\qquad
  \dim\mathcal H_{C(\mathbb F_5)}^{\rm geom}=5^4.
\]

\subsection{End Summary}

The actual Hilbert space in this kinematics is determined by tangent dimension,
not by the residue field alone:
\[
\begin{array}{c|c|c}
\text{local tangent dimension }m & \mathcal H_P^{\rm geom}
  & \text{qudit reading over }K\\ \hline
0 & \mathbb C & \text{no local qudit}\\
1 & \ell^2(K) & \text{one }|K|\text{-level qudit}\\
2 & \ell^2(K^2) & \text{two }|K|\text{-level qudits after a basis choice}.
\end{array}
\]
Dynamics remains separate.  The de Rham CSS construction may later suggest
constraints on these kinematical spaces, but this shard does not impose them.
